\section{The SVD as an Oracle}
\label{sec:svd-oracle}

This section is the reason to read generated register accessors before writing a driver. In this
project it turned two dead ends into two solved ones, in minutes, after days of guessing.

\subsection{Case 1: The Accelerator That Never Ran}

The matrix-vector processor (MVP) was configured exactly as the public material suggested: arrays
filled, dimensions written, one loop, one instruction, then \codeinline{CMD = START}. It never ran.
The result array stayed zero, the status never returned to idle, and nothing faulted — the most
expensive kind of failure, because there is nothing to read.

Two things came from the SVD, in this order:

\begin{enumerate}
    \item \textbf{The block has its own enable register.} The generated code contains an
          \codeinline{EN} register at offset \codeinline{0x04}, described as \emph{Block Enable}, with
          bit~1 as \codeinline{DISABLING}. The firmware had enabled the bus clock
          (\codeinline{CMU\_CLKEN1.MVP}) and never set \codeinline{EN}. A clock gate is not an enable:
          with the block disabled, configuration writes are dropped and the command goes nowhere.
    \item \textbf{The command register's INIT is a qualifier.} The field description says
          ``Initialization Command/Qualifier'', which means it modifies the command written in the
          same access — not a separate initialisation pass, which is what the firmware assumed it had
          to do.
\end{enumerate}

Neither fact is in the tutorial material this project started from. Both are in the file the vendor's
debugger uses.

\subsection{Case 2: A Register That Is Not There}

The GPIO modes taught the same lesson in the other direction. The generated \codeinline{portb\_model}
and \codeinline{porta\_model} are \emph{different types}:

\begin{minted}{rust}
// PORTA and PORTC: sixteen pins, so MODEL and MODEH.
// PORTB and PORTD: eight pins on this package, so MODEL only.
// error[E0609]: no field `portb_modeh` on type `&pac::gpio_s::RegisterBlock`
\end{minted}

A hand-written header would have declared \codeinline{MODEH} for all four ports, because that is what
the family's documentation describes. The SVD describes \emph{this part}, and the compile error it
produces is a fact about the package rather than a limitation of the tool.

\subsection{The Habit}

\begin{enumerate}
    \item Before writing a driver, read the generated module for the peripheral. The doc comments
          carry the offsets and the field names as the silicon team wrote them.
    \item Grep the generated code for the register you think you need. If it is not there, the
          peripheral does not have it on this part — not ``the PAC is incomplete''.
    \item When behaviour disagrees with expectation, read the \emph{neighbouring} registers. In the
          MVP case the answer was four bytes below \codeinline{CFG}: a register that existed, had a
          name, and was never written.
    \item Write what you learn into the crate's docs, not into a chat window: the next person reads
          the accessor's doc comment, not your transcript.
\end{enumerate}
